(1) 以 \(\kappa=k-k_0\) 表示偏离载波的波数，初始高斯包络的傅里叶谱正比 \(\widetilde A(\kappa,0)=A_0s\sqrt{2\pi}\exp(-s^2\kappa^2/2)\)。时间演化中每个谱分量乘 \(\exp[-i(v_g\kappa+\beta\kappa^2/2)t]\)，因此包络为其逆变换。纯二次色散使谱分量幅值不变，但各波数的相位延迟不再线性相同。
(2) 对复高斯积分配方，设 \(D(t)=s^2+i\beta t\)，可得 \(A(x,t)=A_0s/\sqrt{D(t)}\exp[-(x-v_gt)^2/(2D(t))]\)，已略去整体载波相位。强度中心满足 \(x_c(t)=v_gt\)，因为频谱关于 \(\kappa=0\) 对称，二次色散不会额外移动平均群中心。复平方根选择与 \(t=0\) 连续的分支。
(3) 写 \(\xi=x-v_gt\)、\(b=\beta t/s^2\)，强度为 \(|A|^2=|A_0|^2(1+b^2)^{-1/2}\exp[-\xi^2/(s^2(1+b^2))]\)。故其标准差 \(\Delta x(t)=s\sqrt{1+b^2}/\sqrt2\)。复指数相位中出现 \(+\xi^2\beta t/[2(s^4+\beta^2t^2)]\)，意味着不同空间位置的局部波数偏移随 \(\xi\) 线性变化，即色散产生的线性啁啾。
(4) 初始谱强度为 \(|\widetilde A|^2\propto e^{-s^2\kappa^2}\)，因此 \(\Delta k=1/(\sqrt2s)\)。同时初始 \(\Delta x=s/\sqrt2\)，所以 \(\Delta x(0)\Delta k=1/2\)，正好达到高斯傅里叶下界。若 \(\beta=0\)，复宽度 \(D=s^2\) 恒定，整包络只以速度 \(v_g\) 迁移而不展宽，任何后续展宽都来自题设的非零群速度色散。
